% Einheit 2 von 3 – Potenzgesetze (bank/reelle-zahlen/e2.jsonl, zusammenbau v0.1)
%% TODO Zweigzeile Teil 1: was hier gelernt wird (Satz in Schülersprache aus dem Einheitstitel) – Einheit 2
\einheitenkopf[e2]{Einheit 2 von 3 · Potenzgesetze}
\zweigzeile{ab Kl. 9, je nach Buch bis Kl. 10 · keine P10-Aufgabe}
\setcounter{aufgabe}{13}

%% TODO Titel: Ich-kann-Satz fehlt in der Bank (Platzhalter: Kettenname) – Nr. 14
%% TODO Anweisung: ein Satz über der ersten Teilaufgabe fehlt in der Bank – Nr. 14
\begin{aufgabe}{Potenzgesetze}
\begin{teile}
\teil $7^3 \cdot 7^5$ – welches Gesetz passt? \\ \kreuz{gleiche Basis} \\ \kreuz{gleicher Exponent} \\ \kreuz{keins: ausrechnen}
\end{teile}
%% TODO Darstellung für schwach fehlt in der Bank (reelle-zahlen-e2-k1-s1-v1; Abschnitt „Für schwache Schüler“)
\swz{$3^2 \cdot 3^3$ – als Malkette, wie viele Faktoren, als Potenz?}{}
%% TODO Darstellung für schwach fehlt in der Bank (reelle-zahlen-e2-k1-s1-v2; Abschnitt „Für schwache Schüler“)
\swz{$4^2 \cdot 4^2$ – als Malkette, wie viele Faktoren, als Potenz?}{}
%% TODO Darstellung für schwach fehlt in der Bank (reelle-zahlen-e2-k1-s1-v3; Abschnitt „Für schwache Schüler“)
\swz{$5 \cdot 5^3$ – als Malkette, wie viele Faktoren, als Potenz?}{}
%% TODO Darstellung für schwach fehlt in der Bank (reelle-zahlen-e2-k1-s1-v4; Abschnitt „Für schwache Schüler“)
\swz{$2^2 \cdot 2^5$ – als Malkette, wie viele Faktoren, als Potenz?}{}
%% TODO Darstellung für schwach fehlt in der Bank (reelle-zahlen-e2-k1-s1-v5; Abschnitt „Für schwache Schüler“)
\swz{$7 \cdot 7^2$ – als Malkette, wie viele Faktoren, als Potenz?}{}
\swfrage{Was bleibt gleich, was ändert sich?}
%% TODO Darstellung für schwach fehlt in der Bank (reelle-zahlen-e2-k1-s2-v1; Abschnitt „Für schwache Schüler“)
\swz{$5^4 \cdot 5^3$ – als eine Potenz und als Zahl?}{}
%% TODO Darstellung für schwach fehlt in der Bank (reelle-zahlen-e2-k1-s3-v1; Abschnitt „Für schwache Schüler“)
\swz{$7^8 : 7^3$ – als eine Potenz und als Zahl?}{}
%% TODO Darstellung für schwach fehlt in der Bank (reelle-zahlen-e2-k1-s4-v1; Abschnitt „Für schwache Schüler“)
\swz{$(2^3)^2$ – als eine Potenz und als Zahl?}{}
%% TODO Darstellung für schwach fehlt in der Bank (reelle-zahlen-e2-k1-s5-v1; Abschnitt „Für schwache Schüler“)
\swz{$2^4 \cdot 5^4$ – als eine Potenz und als Zahl?}{}
%% TODO Darstellung für schwach fehlt in der Bank (reelle-zahlen-e2-k1-s6-v1; Abschnitt „Für schwache Schüler“)
\swz{$3^2 \cdot 2^3$ – mit Gesetz oder ausrechnen? Ergebnis?}{}
%% TODO Darstellung für schwach fehlt in der Bank (reelle-zahlen-e2-k1-s7-v1; Abschnitt „Für schwache Schüler“)
\swz{$3^2 \cdot 3^2 \cdot 2$ – Ergebnis?}{}
%% TODO Darstellung für schwach fehlt in der Bank (reelle-zahlen-e2-k1-s8-v1; Abschnitt „Für schwache Schüler“)
\swz{$x^5 \cdot x^2$ – vereinfacht?}{}
%% TODO Darstellung für schwach fehlt in der Bank (reelle-zahlen-e2-k1-s9-v1; Abschnitt „Für schwache Schüler“)
\swz{$2x^3 \cdot 5x^4$ – vereinfacht?}{}
%% TODO Darstellung für schwach fehlt in der Bank (reelle-zahlen-e2-k1-s10-v1; Abschnitt „Für schwache Schüler“)
\swz{$3^4 + 3^4 + 3^4$ – als eine Potenz und als Zahl?}{}
%% TODO Darstellung für schwach fehlt in der Bank (reelle-zahlen-e2-k1-s11-v1; Abschnitt „Für schwache Schüler“)
\swz{$(3xy)^2$ – vereinfacht?}{}
\swz[3]{$\frac{6x^5 \cdot 2x^3}{4x^2}$ – vereinfacht?}{}
\end{aufgabe}

%% TODO Titel: Ich-kann-Satz fehlt in der Bank (Platzhalter: Kettenname) – Nr. 15
%% TODO Anweisung: ein Satz über der ersten Teilaufgabe fehlt in der Bank – Nr. 15
\begin{aufgabe}{Potenzen mit negativem Exponenten in Brüche und zurück (Blatt 0 aus potenzen-wurzeln.md; LISUM-PH)}
%% TODO Darstellung für schwach fehlt in der Bank (reelle-zahlen-e2-k2-s1-v1; Abschnitt „Für schwache Schüler“)
\swz{$5^{-2}$ – als Bruch?}{}
\end{aufgabe}

%% TODO Titel: Ich-kann-Satz fehlt in der Bank (Platzhalter: Kettenname) – Nr. 16
%% TODO Anweisung: ein Satz über der ersten Teilaufgabe fehlt in der Bank – Nr. 16
\begin{aufgabe}{Potenzgesetze – Fehler finden, Begründen, Anwendung}
\swz[3]{Finde den Fehler und rechne richtig: \rechnung{3^2 \cdot 3^5 &= 9^7}}{}
\swfrage{Warum ist $3^2 \cdot 3^4$ nicht $9^6$? Begründe mit der Malkette.}
\swz[3]{Ein Speicherbaustein fasst $2^{30}$ Byte. Ein Laufwerk enthält $2^8$ solcher Bausteine. Wie viele Byte fasst das Laufwerk? Gib eine Zweierpotenz an.}{}
\end{aufgabe}

\uebersichtskasten{Gleiche Basis: Beim Malnehmen werden die Exponenten addiert, beim Teilen subtrahiert – die Basis bleibt: 2³ · 2⁴ = 2⁷ (nicht 4⁷ und nicht 2¹²); 5⁶ : 5² = 5⁴. \\ Potenz einer Potenz: Die Exponenten werden multipliziert: (3²)⁴ = 3⁸. \\ Gleicher Exponent: Die Basen werden malgenommen oder geteilt, der Exponent bleibt: 2⁵ · 3⁵ = 6⁵; 10⁴ : 5⁴ = 2⁴. \\ Kein Gesetz für Plus: Potenzen addiert man nur, wenn sie gleich sind – dann zählt man sie: 2⁵ + 2⁵ = 2 · 2⁵ = 2⁶, aber 2³ + 2⁴ = 8 + 16 = 24.}
